\chapter{Tareas y calendario}\label{cap:tareas}

\begin{objbox}
Organizarás lo pendiente con tareas, subtareas y fechas, y las verás en tu calendario habitual.
\end{objbox}

\figura[0.6]{e25-tareas}{Lista de tareas.}
\begin{leyenda}
\item[\mk{1}] \textbf{Calendario}. Genera una dirección (\texttt{.ics}) para suscribir tus tareas y vencimientos en el calendario del Mac o del móvil.
\item[\mk{2}] \textbf{Nueva tarea}.
\item[\mk{3}] \textbf{Casilla}. Marca la tarea como hecha.
\item[\mk{4}] \textbf{Añadir subtarea}.
\item[\mk{5}] \textbf{Editar}.
\item[\mk{6}] \textbf{Eliminar}. Las subtareas se eliminan con su tarea principal.
\end{leyenda}

\figura[0.45]{e26-tarea-nueva}{Ventana de nueva tarea.}
\begin{leyenda}
\item[\mk{1}] \textbf{Título}: «¿Qué hay que hacer?». Obligatorio.
\item[\mk{2}] \textbf{Añadir tarea}. Además puedes fijar fecha, prioridad (alta, normal o baja), cliente vinculado y repetición.
\end{leyenda}
Las tareas con fecha de hoy o atrasadas aparecen en \ui{Hoy}. Una tarea puede estar vinculada a un cliente, a un documento o a un mensaje, de modo que desde la conversación puedes crear la tarea «enviar planos» y encontrarla luego en la ficha del cliente.
